\documentclass{tietreport}

% Core packages
\usepackage[T1]{fontenc}
\usepackage{graphicx}
\usepackage{amsmath}
\usepackage{booktabs}
\usepackage{tabularx}
\usepackage{array}
\usepackage{url}

% Bibliography configuration
\usepackage[
  date=year,
  style=alphabetic,
  backend=biber,
  sorting=nty,
  maxnames=5,
  minnames=3,
  maxalphanames=4,
  minalphanames=3,
  backref=true,
  doi=true,
  isbn=false,
  url=true,
  eprint=true
]{biblatex}
\addbibresource{references.bib}

% Prevent overfull cells in descriptive tables.
\newcolumntype{Y}{>{\raggedright\arraybackslash}X}

%% ==================================================
%% PROJECT INFORMATION
%% ==================================================

\TitlePageHeader{%
  UCS503P -- Software Engineering Project%
}

\ProjectTitle{%
  Meridian: Latent Cartographer%
}

\ProjectSubTitle{%
  A Local-First Spatial Knowledge Engine and Vector Workspace%
}

\author{%
  \texttt{1024160014} Rahul Dadwal \\
  \texttt{1024160133} Jyotsna Sachdeva \\
  \texttt{1024160130} Bisman Singh Rai%
}

\TitlePageSubText{%
  Group: \texttt{[ENTER GROUP NUMBER]} \\
  Bachelor of Engineering, Computer Science and Engineering%
}

\AdvisorName{%
  Dr. Jeelani Asif%
}

\TitlePageFooterText{{%
    \large Mid-Semester Evaluation \\[0.2cm]
    \bfseries Computer Science and Engineering Department%
  } \\[0.15cm]
  Thapar Institute of Engineering and Technology, Patiala%
}

%% ==================================================

\begin{document}

\maketitle
\tableofcontents

\chapter*{Executive Summary}
\addcontentsline{toc}{chapter}{Executive Summary}

Meridian: Latent Cartographer is a local-first spatial knowledge engine proposed as a
Chromium Manifest V3 extension. It addresses the fragmentation created when students,
researchers, and developers work across many browser tabs, articles, documentation pages,
PDFs, Markdown notes, and code snippets. Rather than presenting saved material as a flat
bookmark list, Meridian is intended to represent each item by its semantic meaning and
place related items on an interactive two-dimensional canvas.

The project is designed around on-device processing. Text is converted into semantic
embeddings by the \texttt{Xenova/all-MiniLM-L6-v2} model through Transformers.js and ONNX
Runtime WebAssembly. Search will compare a locally generated query vector with stored item
and passage vectors. IndexedDB is planned for searchable records, while the Origin Private
File System (OPFS) is planned for larger binary assets. UMAP and a collision-correction pass
will later convert item-level vectors into readable canvas coordinates.

Development is incremental. The completed Milestone~1 Foundation Prototype establishes a
loadable Manifest V3 extension, a React and TypeScript application, a dedicated inference
Web Worker, reproducible local model packaging, and real 384-dimensional embedding
generation without external runtime networking. Webpage capture, persistence, semantic
search, and the spatial canvas remain subsequent milestones. This report therefore separates
implemented results from approved design and future work.

% ===== CHAPTER 1: INTRODUCTION =====

\chapter{Introduction}

\section{Background and Context}

Digital research rarely occurs in a single application or document. A typical session may
include dozens of browser tabs, downloaded papers, technical documentation, local notes,
code fragments, and tables. Browsers make these sources easy to open, but their tab bars and
bookmark managers preserve little of the user's conceptual context. Titles are shortened,
similar pages are difficult to distinguish, and a returning user must often reread several
sources before remembering why each one mattered.

Keyword search only partly addresses this problem. A user may remember an idea but not the
words used by the original author. For example, a query about ``storing browser data without
a server'' should be able to retrieve a source discussing IndexedDB or offline persistence
even if the query and source share few literal terms. Dense sentence embeddings provide a
practical mechanism for this form of semantic retrieval. The MiniLM model selected for
Meridian maps text to a 384-dimensional vector space in which semantically related passages
can be compared using cosine similarity~\cite{minilm,transformersjs}.

Many semantic-search products depend on remote inference. Sending research text and search
queries to a cloud service creates privacy, latency, cost, and availability concerns.
Transformers.js makes it possible to execute compatible models locally in the browser using
ONNX Runtime and WebAssembly~\cite{transformersjs,onnxruntime}. Meridian uses this capability
to keep its core embedding and retrieval workflow on the user's device.

\subsection{Problem Statement}

Students, researchers, and developers experience cognitive overload and context loss when
useful information is spread across flat browser-tab lists, bookmarks, local documents, and
code references. Existing organization methods depend heavily on filenames, URLs, folder
placement, or exact keywords. They do not adequately preserve the semantic relationship
between sources or the passage that made a source relevant.

The project must therefore provide a private and responsive method to capture heterogeneous
research material, extract meaningful structure, index the material semantically, persist it
locally, and support visual as well as textual retrieval. It must operate within Manifest V3
lifecycle restrictions and avoid a persistent server or paid inference API.

\subsection{Intended Users}

\begin{itemize}
  \item \textbf{Students} conducting literature reviews, studying documentation, preparing
    reports, or collecting sources across several subjects.
  \item \textbf{Researchers} organizing papers, articles, notes, and supporting evidence by
    conceptual relationship rather than folder membership alone.
  \item \textbf{Developers} saving documentation, technical articles, debugging references,
    code blocks, and structured tables for later retrieval.
\end{itemize}

\subsection{Proposed User Journey}

The target workflow is as follows:

\begin{enumerate}
  \item The user encounters a useful webpage, PDF, Markdown note, or code reference.
  \item The user explicitly saves or imports the source into Meridian.
  \item Meridian extracts and normalizes the source content.
  \item A local model generates semantic embeddings for the item and, in Deep Focus mode,
    its individual chunks.
  \item Searchable records are stored in IndexedDB and large assets are stored in OPFS.
  \item The item appears as a card on the research canvas.
  \item Related items are placed in nearby semantic neighbourhoods.
  \item The user submits a natural-language query.
  \item Meridian ranks stored items and passages using cosine similarity.
  \item The user inspects the matching passage and reopens the original source.
\end{enumerate}

At the present stage, the local embedding portion of this journey is implemented. The other
steps are planned work and are not reported as completed functionality.

\section{Project Objectives}

\begin{enumerate}
  \item \textbf{Establish local semantic inference.} Build a Manifest V3 extension that
    produces real, normalized, 384-dimensional embeddings through a dedicated WASM Web
    Worker, with remote model loading disabled.

  \item \textbf{Implement reliable research capture.} Extract useful text and available
    structure from active webpages, and later support PDFs, Markdown notes, and explicitly
    entered code snippets.

  \item \textbf{Provide semantic retrieval.} Embed natural-language queries locally and rank
    saved item and chunk vectors using cosine similarity.

  \item \textbf{Create a spatial workspace.} Project item-level embeddings onto a stable
    two-dimensional canvas that supports pan, zoom, selection, pinning, manual placement,
    and source inspection.

  \item \textbf{Preserve privacy and offline availability.} Avoid a project backend and paid
    AI APIs for capture, embedding, storage, and search. Package required model and WASM
    assets with the local extension deliverable.

  \item \textbf{Control storage growth.} Use quantized vectors, separate searchable and binary
    storage, on-demand asset loading, and understandable cleanup rules while protecting
    pinned research.

  \item \textbf{Support reproducible team development.} Pin dependencies and model assets,
    define shared TypeScript contracts, separate module ownership, and require independently
    reviewable feature branches.
\end{enumerate}

\section{Scope and Current Boundary}

The completed scope is the technical foundation: extension loading, toolbar launch, typed
worker communication, local model preparation, WASM inference, a diagnostic interface, and
offline verification. The current interface accepts short text and displays embedding
metadata and values.

Webpage capture, IndexedDB and OPFS persistence, semantic search across stored content,
PDF and Markdown ingestion, UMAP projection, Konva-based canvas rendering, lifecycle cleanup,
and workspace export remain within the approved project scope but are not yet implemented.
OCR for scanned PDFs and full language-specific Abstract Syntax Tree analysis are outside the
initial scope.

% ===== CHAPTER 2: METHODOLOGY & DESIGN =====

\chapter{Methodology and Design}

\section{Development Methodology}

The project follows an incremental, risk-first approach. The highest-risk assumption was
whether a Chromium extension could load a pinned model locally, execute it through WASM in a
worker, and generate meaningful embeddings without runtime network access. Milestone~1 was
therefore implemented before capture, storage, and visual design were expanded.

Each milestone has an observable acceptance condition. Shared interfaces are established
before parallel module work so that extraction, storage, inference, and UI code can be
developed independently. Existing academic repository folders and the instructor's C++
sample are preserved. New product code is isolated under \texttt{code/extension/}.

\section{System Architecture}

\subsection{High-Level Design}

Meridian separates orchestration, extraction, computation, persistence, retrieval, layout,
and presentation.

\begin{enumerate}
  \item A Manifest V3 background service worker responds to extension events and coordinates
    browser actions.
  \item Content scripts will extract information from explicitly selected webpages.
  \item A React application will present the capture state, search controls, inspector, and
    spatial workspace.
  \item Dedicated Web Workers perform embedding and later layout computation away from the
    main interface thread.
  \item IndexedDB will hold metadata, text chunks, embeddings, jobs, layout, and settings.
  \item OPFS will hold screenshots and other large binary assets.
\end{enumerate}

The service worker is intentionally not treated as a persistent computation host. Manifest
V3 service workers can be suspended during inactivity, so durable state will be stored and
longer computation will run from an extension page or an appropriate offscreen context. This
design follows the event-driven service-worker model required by Chromium extensions
~\cite{chromeserviceworker,chromeoffscreen}.

\subsection{Implemented Foundation Flow}

The implemented runtime sequence is:

\begin{enumerate}
  \item The user clicks the Meridian toolbar action.
  \item The background service worker uses \texttt{chrome.tabs.create} and an internal
    \texttt{chrome.runtime.getURL} URL to open the application page.
  \item React collects a non-empty short-text input.
  \item The embedding client creates a unique request identifier and posts a typed request to
    a dedicated Web Worker.
  \item The worker loads the packaged tokenizer, quantized ONNX model, and compatible WASM
    runtime from the extension origin.
  \item Transformers.js executes feature extraction with mean pooling and normalization.
  \item The worker verifies that the result contains exactly 384 finite values.
  \item The \texttt{Float32Array} buffer is transferred back to the client.
  \item The diagnostic interface displays the model, dimensions, backend, elapsed time,
    vector norm, and complete vector values.
\end{enumerate}

Inference requests are serialized so that multiple UI requests do not compete for the same
WASM session. The initialized model remains available while the application page and worker
remain open.

\subsection{Planned End-to-End Flow}

The complete processing path will be:

\begin{enumerate}
  \item capture or import a research source;
  \item extract and normalize its content;
  \item create a lightweight representation or Deep Focus chunks;
  \item generate local embeddings;
  \item quantize and persist vectors and metadata;
  \item update semantic coordinates;
  \item render searchable cards; and
  \item retrieve ranked items and passages for a local query.
\end{enumerate}

\subsection{Technical Stack}

\begin{table}[htbp]
\centering
\small
\begin{tabularx}{\textwidth}{@{}p{0.25\textwidth}YY@{}}
\toprule
\textbf{Area} & \textbf{Technology} & \textbf{Role} \\
\midrule
Platform & Chromium Manifest V3 & Extension lifecycle, toolbar action, browser integration,
and content-script execution. \\
Language & TypeScript & Shared contracts and safer communication across extension contexts. \\
Interface & React 19 & Current diagnostic UI and planned workspace controls. \\
Build & Vite 8 & Production bundling for the application, service worker, and inference worker. \\
Local ML & Transformers.js 3.8.1 & Browser-side model loading, tokenization, pooling, and inference. \\
Model & Xenova/all-MiniLM-L6-v2 & Normalized 384-dimensional semantic embeddings. \\
Runtime & ONNX Runtime Web and WASM & Local CPU execution of the quantized ONNX model. \\
Search & Cosine similarity & Planned ranking of query, item, and chunk vectors. \\
Searchable storage & IndexedDB & Planned items, chunks, vectors, jobs, layout, and settings. \\
Binary storage & OPFS & Planned WebP previews and optional snapshots. \\
Extraction & Mozilla Readability & Planned readable webpage extraction with a fallback path. \\
PDF processing & PDF.js & Planned local PDF text extraction and page references. \\
Spatial layout & UMAP-JS & Planned conversion from 384 dimensions to 2D coordinates. \\
Canvas & Konva.js & Planned interactive card rendering, pan, zoom, and selection. \\
Browser automation & Playwright Chromium & Isolated foundation verification in a fresh profile. \\
\bottomrule
\end{tabularx}
\caption{Meridian technology stack and responsibilities}
\label{tab:stack}
\end{table}

\section{Implementation Details}

\subsection{Manifest and Background Module}

The manifest declares a module-based background service worker, toolbar action, minimum
Chrome version 116, and a restrictive Content Security Policy. The current policy permits
scripts, workers, styles, images, and connections only from the extension itself, with the
WASM evaluation capability required by ONNX Runtime. No active-page capture permission has
been added during the foundation milestone because capture is not yet implemented.

The background module currently performs one responsibility: opening the Meridian application
page when the toolbar action is clicked. Keeping it small reduces dependence on transient
service-worker memory and leaves persistent jobs for the later storage layer.

\subsection{Local Embedding Module}

The embedding module uses \texttt{Xenova/all-MiniLM-L6-v2}, a browser-compatible conversion
of the MiniLM sentence-transformer model~\cite{minilm,xenovaminilm}. It requests the quantized
ONNX weights and explicitly selects the WASM backend. One WASM thread is used for compatibility
without cross-origin isolation, and ONNX proxy mode is disabled because inference already
runs in a dedicated worker.

The worker configures the Transformers.js environment to permit local files and disallow
remote models. Browser and filesystem model caches are disabled for the foundation check so
that a warm external cache cannot conceal a missing packaged asset. A worker-level fetch guard
rejects any request whose origin differs from the installed extension origin.

Mean pooling produces a single vector for the input text, and normalization gives the vector
an approximately unit Euclidean norm. For normalized vectors, cosine similarity can be
calculated efficiently using a dot product:

\begin{equation}
  \operatorname{cosine}(\mathbf{u},\mathbf{v}) =
  \frac{\mathbf{u}\cdot\mathbf{v}}
       {\lVert\mathbf{u}\rVert_2\lVert\mathbf{v}\rVert_2}.
\end{equation}

The current UI restricts input to 12,000 characters as a guard, but this does not mean the
entire input always fits the model context. Proper token-aware chunking is scheduled for the
Deep Focus milestone.

\subsection{Reproducible Asset Preparation}

The model is pinned to revision
\texttt{751bff37182d3f1213fa05d7196b954e230abad9}. The project records SHA-256 checksums for
the model configuration, tokenizer files, vocabulary, special-token mapping, and quantized
ONNX model. The preparation script downloads these public assets during development, writes
them atomically, and copies the matching ONNX Runtime WASM glue and binary from the locked
dependency. The production build verifies the checksums and runtime match before bundling.

Downloaded models, dependency folders, build output, caches, and generated verification
results are ignored by Git. A teammate reproduces them from \texttt{package-lock.json},
\texttt{assets.lock.json}, and the preparation script rather than committing large generated
files.

The model and Transformers.js use the Apache License 2.0, and ONNX Runtime uses the MIT
License. Corresponding attribution files are placed beside prepared assets.

\subsection{Shared Contracts and Modularity}

The foundation defines common TypeScript contracts rather than allowing modules to invent
incompatible records.

\begin{table}[htbp]
\centering
\small
\begin{tabularx}{\textwidth}{@{}p{0.30\textwidth}Y@{}}
\toprule
\textbf{Contract} & \textbf{Purpose} \\
\midrule
\texttt{ExtractedDocument} & Represents a webpage, PDF, Markdown note, or snippet with title,
source, normalized text, and ordered paragraph, heading, code, or table blocks. \\
\texttt{CanvasItem} & Represents a workspace card with source metadata, excerpt, canvas
position, pin state, and processing status. \\
\texttt{EmbeddingRequest} & Carries request identity and input text from a client to the worker. \\
\texttt{EmbeddingResponse} & Represents progress, successful result, or failure. \\
\texttt{EmbeddingResult} & Carries the model identity, 384 dimensions, WASM backend,
normalized vector, and elapsed time. \\
\bottomrule
\end{tabularx}
\caption{Shared module contracts}
\label{tab:contracts}
\end{table}

Module ownership is separated into background, content, workers, extraction, storage, search,
layout, UI, and shared contracts. Extraction work can therefore return
\texttt{ExtractedDocument} records while UI work uses \texttt{CanvasItem} records. Changes to
the manifest, package files, build configuration, embedding worker, or shared contracts require
coordination because these files affect multiple modules.

\subsection{Planned Capture and Extraction Module}

Active-page capture will be explicit. Mozilla Readability will isolate the main article where
possible, while a safe DOM-text fallback will handle pages for which article extraction fails
~\cite{readability}. Structured extraction will preserve block order, headings, paragraphs,
\texttt{pre/code} blocks, and tables. A repeated source will offer an update instead of
silently creating a duplicate.

Lightweight mode will retain representative text, metadata, one available preview, and one
item vector. Deep Focus mode will retain full text and create proposed 256-token windows with
a 50-token overlap. Actual chunk boundaries will use the model tokenizer and account for
special tokens and input limits.

Local PDFs will later be processed through PDF.js~\cite{pdfjs}. Markdown import will preserve
headings, code fences, tables, and note text. Scanned PDFs without an extractable text layer
will receive an explanatory message; OCR is outside the initial scope.

\subsection{Planned Storage Module}

IndexedDB will act as the hot searchable layer for item metadata, chunks, vectors,
quantization parameters, model identifiers, processing jobs, coordinates, and settings
~\cite{indexeddb}. OPFS will hold larger screenshots and optional snapshots that are loaded on
demand~\cite{opfs}.

The planned job states are \texttt{queued}, \texttt{extracting}, \texttt{embedding},
\texttt{ready}, and \texttt{failed}. Persisting these states allows recovery after a service
worker or application interruption and prevents unsafe duplicate retries.

SQ8 quantization is planned for stored embeddings. The implementation must preserve the scale
or range values required to reconstruct similarity scores and must compare retrieval quality
with the original floating-point vectors before adopting compression. The project target is
at least 70\% vector-storage reduction, but this target has not yet been measured.

\subsection{Planned Search and Canvas Modules}

The search module will embed the query using the same local model and rank stored item and
chunk vectors with cosine similarity. Search ranking will use high-dimensional vectors rather
than distance on the 2D map. The selected result will show its relevant passage and source.

The layout module will use UMAP to project item-level vectors to two dimensions
~\cite{umap,umapjs}. A collision-correction pass will reduce card overlap. Coordinates will be
persisted, manually placed or pinned cards will be respected, and users will request major
reorganization explicitly. Konva.js is planned for the interactive canvas~\cite{konva}.
Viewport culling and controlled image loading will reduce rendering and memory pressure.

\subsection{Storage Lifecycle and Portability}

Pinned items will be protected from automatic cleanup. Unpinned preview assets are proposed
to become eligible for cleanup after seven inactive days while text and searchable vectors
remain. Thirty-day centroid-only reduction is optional and disabled initially because it
would reduce passage-level search quality.

The planned portability layer includes versioned JSON export and import with preservation of
content and coordinates. A standalone HTML export will initially act as an inert viewer.
Including local AI search inside that exported file requires a separate model-distribution
decision, and captured scripts must never execute in the viewer.

% ===== CHAPTER 3: RESULTS & EVALUATION =====

\chapter{Results and Evaluation}

\section{Evaluation Methodology}

Milestone~1 was evaluated as an actual extension rather than only as isolated JavaScript.
TypeScript compilation and the Vite production build first confirmed that the application,
background worker, and inference worker bundled successfully. The build also verified model
checksums and compatibility of the copied WASM runtime.

An integration scenario then loaded the built extension into a fresh Playwright Chromium
profile. Before the application page was opened, browser offline mode was enabled. External
browser traffic was additionally directed to an unreachable local proxy, DNS resolution was
disabled, and background networking was suppressed. This provided multiple independent
controls against accidental external access.

Three short sentences were embedded:

\begin{enumerate}
  \item ``A cat is sleeping on a warm windowsill.''
  \item ``A kitten is resting beside the sunny window.''
  \item ``Database transactions provide atomicity and isolation.''
\end{enumerate}

The evaluation checked that each result contained 384 finite values, was approximately unit
normalized, contained non-constant values, and differed across unrelated inputs. It also
confirmed that related sentences received greater cosine similarity than an unrelated pair.
Every model and runtime request observed during inference used a
\texttt{chrome-extension://} URL. A separate external navigation probe failed under the same
browser restrictions.

No student usability study has yet been conducted. The proposed pilot involving engineering
students remains future work.

\section{Foundation Results}

\begin{table}[htbp]
\centering
\small
\begin{tabularx}{\textwidth}{@{}Yp{0.27\textwidth}p{0.23\textwidth}@{}}
\toprule
\textbf{Measure} & \textbf{Project target or rule} & \textbf{Observed foundation result} \\
\midrule
Embedding dimensions & Exactly 384 & 384 finite values \\
Embedding norm & Approximately 1 after normalization & 1.000000 in the diagnostic result \\
Runtime backend & Local WASM & WASM \\
First short embedding & No final requirement established & Approximately 360.2 ms \\
Later short embeddings & No final requirement established & Approximately 7.7 ms and 7.6 ms \\
Related-sentence similarity & Greater than unrelated comparison & Approximately 0.589 \\
Unrelated-sentence similarity & Lower than related comparison & Approximately $-0.069$ \\
External inference requests & None & None observed \\
External navigation probe & Must fail in offline scenario & Blocked \\
Semantic search latency & Median $\leq 20$ ms over 200+ artifacts & Not yet measurable \\
Canvas smoothness & 55--60 FPS with 100+ active nodes & Not yet measurable \\
Tab ingestion & $\leq 1.5$ s per tab & Not yet measurable \\
Vector storage reduction & At least 70\% compared with Float32 & Not yet measurable \\
\bottomrule
\end{tabularx}
\caption{Foundation observations and outstanding project targets}
\label{tab:results}
\end{table}

The first-call time includes model and WASM initialization. Later calls reuse the model held
by the worker and should not be interpreted as complete search latency. The proposed 20 ms
retrieval target includes a different workload and will be measured after query embedding,
stored-vector scoring, result selection, and UI response have been integrated.

\section{Analysis}

\subsection{Objectives Met in the Current Milestone}

The foundation objective was met. A real quantized MiniLM model generated normalized
384-dimensional vectors inside a dedicated extension worker. The extension opened through a
Manifest V3 background action, and the model and WASM assets were served from the installed
extension during the offline scenario.

Semantic comparison behaved in the expected direction: the cat and kitten sentences were
substantially closer than the cat and database sentences. This does not constitute a complete
retrieval-quality study, but it rules out a constant or obviously unrelated mock result and
supports continuation with the chosen embedding pipeline.

\subsection{Engineering Challenges and Resolutions}

\begin{itemize}
  \item \textbf{Offline asset loading.} Transformers.js normally supports hosted model and
    runtime paths. Meridian instead pins and packages the model and compatible WASM files,
    disables remote models, verifies checksums, and restricts worker fetches.

  \item \textbf{Manifest V3 lifecycle.} Persistent model state is not placed in the background
    service worker. The current application page owns the inference worker, while later job
    state will be persisted so that browser suspension does not lose progress.

  \item \textbf{Reliable verification.} A fresh profile, disabled model caches, browser offline
    mode, an unreachable proxy, disabled DNS, observed request URLs, and an external negative
    control reduce the possibility that cached or remote files caused an apparent pass.

  \item \textbf{Dependency advisory.} The Node-side \texttt{sharp} dependency inherited by
    Transformers.js was overridden to version 0.35.0. The dependency audit subsequently
    reported no known vulnerabilities. The extension's runtime path uses the browser WASM
    bundle rather than Sharp or ONNX Runtime Node.
\end{itemize}

\subsection{Current Limitations}

\begin{itemize}
  \item The current UI is an embedding diagnostic, not the complete Meridian workspace.
  \item Webpage capture, structured extraction, duplicate handling, and screenshots are not
    implemented.
  \item IndexedDB and OPFS schemas are not implemented, so research items do not yet persist.
  \item Semantic retrieval across saved content and passage inspection are not implemented.
  \item The spatial canvas, UMAP projection, collision correction, and viewport culling are
    not implemented.
  \item PDF, Markdown, and snippet ingestion are not implemented.
  \item The current input guard is character-based; token-aware document chunking remains
    future work.
  \item The extension is locally deployed as an unpacked build and is not published in the
    Chrome Web Store.
  \item Reported times are observations from one environment, not product-wide guarantees.
\end{itemize}

% ===== CHAPTER 4: CONCLUSION =====

\chapter{Conclusion}

\section{Summary of Work}

The Meridian project has progressed from an architectural proposal to a working foundation
prototype. The current implementation establishes a Chromium Manifest V3 application using
React, TypeScript, and Vite; a background action that opens the application; a typed client
and dedicated inference worker; pinned and checksum-verified MiniLM assets; local ONNX WASM
execution; and a reproducible offline verification workflow.

The implementation is contained within \texttt{code/extension/}, preserving the instructor's
repository structure, academic reports, journals, and existing C++ example. The foundation
is recorded in commit \texttt{1253aae} on the remote \texttt{codex/foundation} branch. At the
latest recorded repository review, this branch had not yet been merged into the default
\texttt{master} branch.

\section{Key Findings}

\begin{itemize}
  \item \textbf{Local browser inference is feasible.} A packaged quantized MiniLM model can
    generate real 384-dimensional semantic embeddings inside a Manifest V3 extension using
    WASM and without runtime cloud inference.

  \item \textbf{Worker separation supports the intended architecture.} Model computation can
    remain outside the React UI and transient background service worker while still exposing
    a small typed client interface to future search and ingestion modules.

  \item \textbf{Offline claims require explicit controls.} Packaging files alone is not enough.
    Remote-model configuration, extension CSP, origin-restricted fetches, fresh-profile
    execution, cache control, and negative network checks together provide stronger evidence.

  \item \textbf{Measured foundation timings do not prove final search targets.} Initialization,
    query embedding, scoring, UI response, ingestion, and layout are different measurements
    and must be reported separately in later milestones.
\end{itemize}

\section{Future Work and Recommendations}

\begin{enumerate}
  \item \textbf{Merge and stabilize the foundation.} Review the
    \texttt{codex/foundation} branch and merge it into \texttt{master} before parallel feature
    branches begin.

  \item \textbf{Implement capture and persistence.} Add explicit active-page capture,
    Readability extraction, structured blocks, duplicate handling, an IndexedDB schema,
    initial OPFS assets, persistent job states, and browser-restart recovery.

  \item \textbf{Deliver the first complete workflow.} Connect stored items to the embedding
    worker, cosine-similarity search, a basic canvas, an inspector, result highlighting, and
    source reopening.

  \item \textbf{Add Deep Focus ingestion.} Implement tokenizer-aware overlapping chunks,
    PDF and Markdown import, manual snippets, source references, and passage-level retrieval.

  \item \textbf{Build the semantic map.} Add UMAP projection, collision correction, stable
    coordinates, pinned and manually positioned cards, explicit reorganization, and viewport
    culling.

  \item \textbf{Complete storage and portability.} Add understandable cleanup controls,
    pinned-item protection, versioned JSON export and import, and a safe standalone HTML
    viewer.

  \item \textbf{Evaluate at project scale.} Measure retrieval quality and latency on at least
    200 artifacts, canvas responsiveness with at least 100 nodes, actual storage reduction,
    cold and warm inference, ingestion time, and usability with an agreed pilot protocol.

  \item \textbf{Prepare distribution.} Produce a reproducible extension package and complete
    project documentation before considering Chrome Web Store publication.
\end{enumerate}

\section{Final Remarks}

Milestone~1 does not yet solve the entire research-organization problem, but it validates the
most important technical dependency beneath the proposed solution: meaningful local semantic
inference inside the browser. The next phase can now build capture and storage against a real
embedding service rather than a mock. Maintaining the distinction between implemented,
verified, approved, and planned capabilities will keep the project report accurate as Meridian
progresses toward its spatial research workspace.

% ===== BIBLIOGRAPHY =====

\printbibliography[heading=bibintoc]

\end{document}
